% subsection content: 等差型 \, $a_{n+1}=a_n+d$

まずはもっともシンプルな漸化式であろう,等差型について説明していく.
この漸化式は$a_1$に対して$d$がどんどんと足されていく漸化式である.
一般項は簡単な代入計算によって容易に示せる.
実際にやってみると
\begin{align*}
   & a_{n+1}=a_n+d                   \\
   & a_{n+1}=(a_{n-1}+d)+d           \\
   & a_{n+1}=\{(a_{n-2}+d)+d\}+d     \\
   & \qquad \vdots                   \\
   & a_{n+1}=a_1+d+d+\cdots+d        \\
   & a_{n+1}=a_1+nd                  \\
   & \therefore \quad a_n=a_1+(n-1)d
\end{align*}
となり一般項が示される.
1つ前の項をどんどん代入していくのだ.
数式の3行目から5行目周辺の式変形がすこし心配かもしれないが,\,$d$の個数と$a$の添字部分の和は常に$n+1$となることから考えれば分かるだろう.
実際,\,$a_{n+1}$に至るまでに何度$d$が足されたかを考えているわけだから,心配であれば実際に適当な$n$で何度か実験してみるのもいい手だろう.

実際に例題を与えよう.

\noindent \textbf{【例題】}
\begin{enumerate}[label=(\arabic*),parsep=3pt]
  \item $\displaystyle a_1=1,\, a_{n+1}=a_n+1$
  \item $\displaystyle a_1=\frac{1}{2},\, a_{n+1}=a_n+\frac{1}{4}$
  \item $\displaystyle a_1=-\sqrt{3},\, a_{n+1}=a_n-\sqrt{2}$
\end{enumerate}

\noindent\textbf{【方針】}
\begin{enumerate}[label=(\arabic*),parsep=3pt]
  \item 等差型の一般項
        $\displaystyle a_n=a_1+(n-1)d$
        を用います.
  \item 公差は$\displaystyle d=\frac14$です.
        等差型の一般項に代入します.
  \item 公差は$\displaystyle d=-\sqrt2$です.
        負の公差の場合も同じ公式を使います.
\end{enumerate}

\noindent\textbf{【解答】}
\begin{enumerate}[label=(\arabic*),parsep=3pt]
  \item $\displaystyle a_n=1+(n-1)=n$
  \item $\displaystyle a_n=\frac12+\frac{n-1}{4}
          =\frac{n+1}{4}$
  \item $\displaystyle a_n=-\sqrt3-(n-1)\sqrt2$
\end{enumerate}

\noindent\textbf{【解法】}\\
解法略.
